\documentclass[11pt]{article}
\usepackage[margin=1in]{geometry}
\usepackage[T1]{fontenc}
\usepackage{lmodern}
\usepackage{microtype}
\usepackage[table]{xcolor}
\usepackage{amsmath,amssymb,amsfonts,mathtools,bm}
\usepackage{physics}
\usepackage{booktabs,longtable,array}
\usepackage{enumitem}
\usepackage[hidelinks]{hyperref}
\usepackage[most]{tcolorbox}
\usepackage{fancyhdr}
\usepackage{lastpage}
\usepackage{titlesec}
\usepackage{mathrsfs}

\definecolor{Navy}{HTML}{1F3A5F}
\definecolor{SoftBlue}{HTML}{EAF2F8}
\definecolor{SoftGray}{HTML}{F7F9FA}
\definecolor{RuleGray}{HTML}{D6DBDF}

\pagestyle{fancy}
\fancyhf{}
\lhead{MHD\_PINN\_proj1}
\rhead{Module 1 Physics Notes}
\cfoot{\thepage\ of \pageref{LastPage}}
\setlength{\headheight}{14pt}
\renewcommand{\headrulewidth}{0.4pt}
\renewcommand{\footrulewidth}{0pt}

\titleformat{\section}{\large\bfseries\color{Navy}}{\thesection}{0.6em}{}
\titleformat{\subsection}{\normalsize\bfseries\color{Navy}}{\thesubsection}{0.6em}{}
\titleformat{\subsubsection}{\normalsize\bfseries}{\thesubsubsection}{0.6em}{}

\newcommand{\uvec}{\mathbf{u}}
\newcommand{\Bvec}{\mathbf{B}}
\newcommand{\Jvec}{\mathbf{J}}
\newcommand{\Uvec}{\mathbf{U}}
\newcommand{\Fvec}{\mathbf{F}}
\newcommand{\Gvec}{\mathbf{G}}
\newcommand{\Iten}{\mathbf{I}}
\newcommand{\ptot}{p_{\mathrm{tot}}}

\newtcolorbox{overviewbox}{
  colback=SoftBlue,
  colframe=Navy,
  boxrule=0.7pt,
  arc=2pt,
  left=8pt,right=8pt,top=8pt,bottom=8pt
}

\newtcolorbox{physicsbox}{
  colback=SoftGray,
  colframe=RuleGray,
  boxrule=0.5pt,
  arc=2pt,
  left=8pt,right=8pt,top=8pt,bottom=8pt
}

\begin{document}

\begin{center}
    {\LARGE\bfseries Module 1: Ideal MHD Equations and Detailed Derivations}\\[0.6em]
    {\large MHD\_PINN\_proj1 Physics Notes}\\[0.4em]
    {\normalsize Revised to include step-by-step derivations of the conservative equations}
\end{center}

\vspace{0.8em}

\begin{overviewbox}
These notes derive the ideal magnetohydrodynamics (MHD) equations used by the project solver from basic conservation laws and electromagnetic force balance. The goal is not just to state the final equations, but to show how each one is obtained. In particular, the derivation of the momentum equation is carried out in detail from Newton's second law, through the Euler equation of compressible gas dynamics, through the Lorentz force density, and finally into the magnetic-stress-tensor form used by the finite-volume code.
\end{overviewbox}

\section{Symbol and variable table}
\rowcolors{2}{SoftGray}{white}
\begin{longtable}{>{\raggedright\arraybackslash}p{0.18\textwidth} >{\raggedright\arraybackslash}p{0.28\textwidth} >{\raggedright\arraybackslash}p{0.46\textwidth}}
\toprule
\textbf{Symbol} & \textbf{Name} & \textbf{Meaning} \\
\midrule
\endfirsthead
\toprule
\textbf{Symbol} & \textbf{Name} & \textbf{Meaning} \\
\midrule
\endhead
$\rho$ & mass density & Mass per unit volume of the plasma. \\
$p$ & gas pressure & Thermal pressure of the fluid. \\
$\gamma$ & ratio of specific heats & Closure parameter for an ideal gas, commonly $\gamma=5/3$. \\
$\uvec=(u_x,u_y,u_z)$ & fluid velocity & Plasma bulk-flow velocity. \\
$\Bvec=(B_x,B_y,B_z)$ & magnetic field & Magnetic-field vector. \\
$\Jvec$ & current density & Electric-current density carried by the plasma. \\
$E$ & total energy density & Sum of internal, kinetic, and magnetic energy densities. \\
$\ptot$ & total pressure & Sum of gas pressure and magnetic pressure. \\
$\Uvec$ & conserved-state vector & Vector of conserved variables updated by the finite-volume solver. \\
$\Fvec$ & $x$-direction flux vector & Flux of conserved quantities through a face normal to the $x$-direction. \\
$\Gvec$ & $y$-direction flux vector & Flux of conserved quantities through a face normal to the $y$-direction. \\
$\Iten$ & identity tensor & Tensor that leaves vectors unchanged under multiplication. \\
$\uvec\uvec$ & dyadic product of velocity & Rank-2 tensor with components $(u_i u_j)$. \\
$\Bvec\Bvec$ & dyadic product of field & Rank-2 tensor with components $(B_i B_j)$. \\
$\uvec\cdot\Bvec$ & velocity--field dot product & Measures how much the flow is aligned with the magnetic field. \\
$\nabla\cdot\Bvec$ & magnetic divergence & Must be zero in ideal MHD because magnetic monopoles are not observed. \\
$x,y,t$ & coordinates and time & Two spatial coordinates and time variable used by the solver. \\
\bottomrule
\end{longtable}

\section{Why MHD is the model used by the project}
A plasma contains charged particles, so it feels both ordinary fluid forces and electromagnetic forces. If one tried to track every particle separately, the description would become microscopic and extremely expensive. MHD instead describes the plasma in terms of coarse-grained fields such as density, velocity, pressure, and magnetic field. The basic physical picture is:
\begin{itemize}[leftmargin=1.8em]
    \item mass is transported by the bulk flow,
    \item momentum changes because of pressure gradients, inertia, and magnetic forces,
    \item energy is carried and converted between internal, kinetic, and magnetic forms,
    \item the magnetic field evolves with the fluid motion in the ideal-MHD limit.
\end{itemize}

\begin{physicsbox}
\textbf{Key idea.} Gas dynamics already has inertia and gas pressure. MHD keeps those but adds magnetic stress. Magnetic stress has two pieces: an isotropic part that behaves like pressure, and an anisotropic part that behaves like tension along field lines.
\end{physicsbox}

\section{The conserved-state vector}
The solver evolves the conserved variable vector
\begin{equation}
\Uvec =
\begin{bmatrix}
\rho \\
\rho u_x \\
\rho u_y \\
\rho u_z \\
E \\
B_x \\
B_y \\
B_z
\end{bmatrix}.
\label{eq:Uvec}
\end{equation}
This choice comes directly from the conservation laws:
\begin{itemize}[leftmargin=1.8em]
    \item the first entry is mass density because mass is conserved,
    \item the next three entries are the three components of momentum density because Newton's second law is a momentum balance law,
    \item the fifth entry is total energy density because the ideal-MHD system conserves total energy,
    \item the final three entries are the magnetic-field components because the induction equation evolves the magnetic field dynamically.
\end{itemize}

\section{Total energy density}
The total energy density is
\begin{equation}
E = \frac{p}{\gamma-1} + \frac{1}{2}\rho |\uvec|^2 + \frac{1}{2}|\Bvec|^2.
\label{eq:total-energy}
\end{equation}
The three terms are:
\begin{enumerate}[leftmargin=1.8em]
    \item \textbf{internal energy density} $e_{\mathrm{int}}=p/(\gamma-1)$ for an ideal gas,
    \item \textbf{kinetic energy density} $\frac12 \rho |\uvec|^2$,
    \item \textbf{magnetic energy density} $\frac12 |\Bvec|^2$ in the normalized units used by the code.
\end{enumerate}
The code stores $E$ instead of storing internal energy separately because the conservative finite-volume method updates one total-energy balance law.

\section{Ideal-MHD equations in compact conservative form}
The full system can be written as
\begin{equation}
\pdv{\Uvec}{t} + \pdv{\Fvec(\Uvec)}{x} + \pdv{\Gvec(\Uvec)}{y} = 0.
\label{eq:cons-law}
\end{equation}
Equation \eqref{eq:cons-law} is compact, but each row corresponds to a physically distinct balance law. The rest of the note derives those rows carefully.

\section{Mass conservation in detail}
Consider a fixed control volume $V$ with boundary surface $\partial V$ and outward unit normal $\hat{\mathbf{n}}$. The total mass inside the control volume is
\begin{equation}
M_V(t)=\int_V \rho\, dV.
\label{eq:mass-control}
\end{equation}
Mass conservation says that the only way the mass inside $V$ can change is by mass crossing the boundary. The mass flux through a unit area is $\rho \uvec\cdot\hat{\mathbf{n}}$, so
\begin{equation}
\dv{}{t}\int_V \rho\, dV = -\int_{\partial V} \rho\uvec\cdot\hat{\mathbf{n}}\, dS.
\label{eq:mass-integral}
\end{equation}
Use the divergence theorem,
\begin{equation}
\int_{\partial V} \rho\uvec\cdot\hat{\mathbf{n}}\, dS = \int_V \nabla\cdot(\rho\uvec)\, dV,
\label{eq:div-thm-mass}
\end{equation}
and write the time derivative inside the volume integral:
\begin{equation}
\int_V \pdv{\rho}{t}\, dV = -\int_V \nabla\cdot(\rho\uvec)\, dV.
\label{eq:mass-volume-form}
\end{equation}
Because this must hold for any control volume $V$, the integrands must match pointwise. Therefore,
\begin{equation}
\pdv{\rho}{t} + \nabla\cdot(\rho\uvec)=0.
\label{eq:mass}
\end{equation}
In two-dimensional Cartesian form, the solver uses
\begin{equation}
\pdv{\rho}{t} + \pdv{}{x}(\rho u_x) + \pdv{}{y}(\rho u_y)=0.
\label{eq:mass2d}
\end{equation}

\section{Momentum conservation and magnetic stress}
This is the most important derivation in the note because it explains exactly where the momentum fluxes in the code come from.

\subsection{Newton's second law for a small fluid element}
Take a small fluid element of volume $\Delta V$ and mass $\rho\,\Delta V$. Newton's second law says
\begin{equation}
(\rho\,\Delta V)\,\mathbf{a} = \mathbf{F}_{\text{net on element}}.
\label{eq:newton-small-element}
\end{equation}
Divide by $\Delta V$:
\begin{equation}
\rho\,\mathbf{a} = \frac{\mathbf{F}_{\text{net on element}}}{\Delta V}.
\label{eq:newton-per-volume}
\end{equation}
So Newton's second law ``per unit volume'' means: acceleration times mass density equals force density.

The acceleration of a moving fluid parcel is not just $\partial \uvec/\partial t$ at one point. A parcel can move into a region where the velocity field differs. Therefore the parcel acceleration is the material derivative,
\begin{equation}
\mathbf{a}=\frac{D\uvec}{Dt}=\pdv{\uvec}{t}+(\uvec\cdot\nabla)\uvec.
\label{eq:material-accel}
\end{equation}
Substituting \eqref{eq:material-accel} into \eqref{eq:newton-per-volume} gives the general momentum balance law
\begin{equation}
\rho\frac{D\uvec}{Dt}=\mathbf{f},
\label{eq:newton-fluid}
\end{equation}
where $\mathbf{f}$ is the total force per unit volume.

\subsection{Pressure force on a fluid element}
Now derive the pressure-force density explicitly. Pressure acts normal to a surface and pushes inward on a fluid element. Consider the $x$-direction. The left face of the element sees pressure $p(x)$ pushing in the $+x$ direction, and the right face sees pressure $p(x+\Delta x)$ pushing in the $-x$ direction. The net pressure force in the $x$ direction is approximately
\begin{equation}
F_x^{(p)} = p(x)\,\Delta y\Delta z - p(x+\Delta x)\,\Delta y\Delta z.
\label{eq:pressure-force-x}
\end{equation}
Expand $p(x+\Delta x)$ to first order,
\begin{equation}
p(x+\Delta x)=p(x)+\pdv{p}{x}\Delta x + \mathcal{O}(\Delta x^2).
\label{eq:taylor-pressure}
\end{equation}
Then
\begin{equation}
F_x^{(p)} \approx -\pdv{p}{x}\Delta x\Delta y\Delta z = -\pdv{p}{x}\Delta V.
\label{eq:pressure-force-x-2}
\end{equation}
Dividing by $\Delta V$ gives the $x$-component of pressure-force density:
\begin{equation}
\frac{F_x^{(p)}}{\Delta V} = -\pdv{p}{x}.
\label{eq:pressure-force-density-x}
\end{equation}
Repeating the same argument in $y$ and $z$ gives the vector form
\begin{equation}
\mathbf{f}_p = -\nabla p.
\label{eq:pressure-force-density}
\end{equation}
Substituting \eqref{eq:pressure-force-density} into \eqref{eq:newton-fluid} for a neutral, inviscid compressible fluid with no magnetic field gives
\begin{equation}
\rho\frac{D\uvec}{Dt} = -\nabla p.
\label{eq:euler-lagrangian}
\end{equation}
This is the Euler equation in nonconservative form.

\subsection{From Newton's law to the conservative Euler momentum equation}
The finite-volume solver does not evolve parcel acceleration directly. It evolves the momentum density $\rho\uvec$. We now show how \eqref{eq:euler-lagrangian} becomes a conservation law.

Start with the product rule for $\rho\uvec$:
\begin{equation}
\pdv{}{t}(\rho\uvec)=\rho\pdv{\uvec}{t}+\uvec\pdv{\rho}{t}.
\label{eq:product-rule-time}
\end{equation}
Now expand the divergence of the dyadic tensor $\rho\uvec\uvec$ in index notation. Its $i$th component is
\begin{equation}
\left[\nabla\cdot(\rho\uvec\uvec)\right]_i = \partial_j(\rho u_i u_j).
\label{eq:rho-uu-div-index-start}
\end{equation}
Apply the product rule:
\begin{equation}
\partial_j(\rho u_i u_j)=u_i\,\partial_j(\rho u_j)+\rho u_j\partial_j u_i.
\label{eq:rho-uu-div-index-expanded}
\end{equation}
Recognizing that $\partial_j(\rho u_j)=\nabla\cdot(\rho\uvec)$ and $u_j\partial_j u_i=[(\uvec\cdot\nabla)\uvec]_i$, this becomes
\begin{equation}
\nabla\cdot(\rho\uvec\uvec)=\uvec\,\nabla\cdot(\rho\uvec)+\rho(\uvec\cdot\nabla)\uvec.
\label{eq:rho-uu-div-vector}
\end{equation}
Add \eqref{eq:product-rule-time} and \eqref{eq:rho-uu-div-vector}:
\begin{equation}
\pdv{}{t}(\rho\uvec)+\nabla\cdot(\rho\uvec\uvec)
=\rho\pdv{\uvec}{t}+\rho(\uvec\cdot\nabla)\uvec+\uvec\left[\pdv{\rho}{t}+\nabla\cdot(\rho\uvec)\right].
\label{eq:momentum-preidentity}
\end{equation}
The bracketed term vanishes by mass conservation \eqref{eq:mass}. Therefore,
\begin{equation}
\pdv{}{t}(\rho\uvec)+\nabla\cdot(\rho\uvec\uvec)=\rho\left(\pdv{\uvec}{t}+(\uvec\cdot\nabla)\uvec\right)=\rho\frac{D\uvec}{Dt}.
\label{eq:mom-identity}
\end{equation}
Substitute \eqref{eq:euler-lagrangian} into \eqref{eq:mom-identity}:
\begin{equation}
\pdv{}{t}(\rho\uvec)+\nabla\cdot(\rho\uvec\uvec) = -\nabla p.
\label{eq:euler-precons}
\end{equation}
Now write the pressure-gradient term as the divergence of an isotropic tensor:
\begin{equation}
-\nabla p = -\nabla\cdot(p\Iten),
\label{eq:pressure-tensor-form}
\end{equation}
because componentwise,
\begin{equation}
\left[\nabla\cdot(p\Iten)\right]_i = \partial_j(p\delta_{ij}) = \partial_i p.
\label{eq:pressure-tensor-proof}
\end{equation}
Therefore the Euler momentum equation becomes
\begin{equation}
\pdv{}{t}(\rho\uvec) + \nabla\cdot\left[\rho\uvec\uvec + p\Iten\right] = 0.
\label{eq:euler-cons}
\end{equation}
This is the full answer to ``how does Eq.~(6) come from Newton's second law?'' The chain is:
\begin{enumerate}[leftmargin=1.8em]
    \item Newton's second law for a fluid parcel,
    \item pressure-force density $-\nabla p$,
    \item material derivative for parcel acceleration,
    \item continuity equation to convert from acceleration form to momentum-conservation form,
    \item tensor form $p\Iten$ to place the pressure term inside a divergence.
\end{enumerate}

\subsection{Add the magnetic force density}
In MHD, the conducting fluid also feels the Lorentz force density
\begin{equation}
\mathbf{f}_{B}=\Jvec\times\Bvec.
\label{eq:lorentz-density}
\end{equation}
Adding this to \eqref{eq:euler-cons} gives
\begin{equation}
\pdv{}{t}(\rho\uvec) + \nabla\cdot\left[\rho\uvec\uvec + p\Iten\right] = \Jvec\times\Bvec.
\label{eq:mhd-mom-with-lorentz}
\end{equation}
So the remaining task is to rewrite the magnetic force as the divergence of a stress tensor.

\subsection{Express the Lorentz force only in terms of $\Bvec$}
In the normalized ideal-MHD units used by the solver,
\begin{equation}
\Jvec = \nabla\times\Bvec.
\label{eq:ampere}
\end{equation}
Therefore,
\begin{equation}
\Jvec\times\Bvec = (\nabla\times\Bvec)\times\Bvec.
\label{eq:lorentz-as-B}
\end{equation}
Use the vector identity
\begin{equation}
(\nabla\times\Bvec)\times\Bvec = (\Bvec\cdot\nabla)\Bvec - \nabla\left(\frac12 |\Bvec|^2\right) + \Bvec(\nabla\cdot\Bvec).
\label{eq:lorentz-identity}
\end{equation}
Ideal MHD imposes
\begin{equation}
\nabla\cdot\Bvec = 0,
\label{eq:divB0}
\end{equation}
so
\begin{equation}
\Jvec\times\Bvec = (\Bvec\cdot\nabla)\Bvec - \nabla\left(\frac12 |\Bvec|^2\right).
\label{eq:lorentz-pressure-tension}
\end{equation}
This already shows the physical split:
\begin{itemize}[leftmargin=1.8em]
    \item $-\nabla(\frac12 |\Bvec|^2)$ is the force caused by spatial variation of magnetic energy density, so it behaves like an isotropic magnetic pressure,
    \item $(\Bvec\cdot\nabla)\Bvec$ acts along the field and depends on changes of field direction and strength along field lines, so it behaves like tension.
\end{itemize}

\begin{physicsbox}
\textbf{Why magnetic pressure is isotropic.}
The scalar $\frac12 |\Bvec|^2$ has no direction attached to it. Its gradient therefore points from stronger-field regions toward weaker-field regions. That is the same mathematical structure as an ordinary pressure gradient, so it contributes an isotropic pressure-like stress.

\textbf{Why magnetic tension is directional.}
The term $(\Bvec\cdot\nabla)\Bvec$ differentiates the field in the direction of the field itself. It therefore measures how the field changes along its own lines, which is why it represents directional line-stretching or line-straightening tension.
\end{physicsbox}

\subsection{Magnetic force as the divergence of a stress tensor}
To combine the magnetic force with the fluid momentum flux, we must write it as a divergence. First,
\begin{equation}
\nabla\cdot\left(\frac12 |\Bvec|^2\Iten\right)=\nabla\left(\frac12 |\Bvec|^2\right),
\label{eq:isotropic-div}
\end{equation}
for exactly the same reason that $\nabla\cdot(p\Iten)=\nabla p$.

Now consider the tensor $\Bvec\Bvec$. Its $i$th divergence component is
\begin{equation}
\left[\nabla\cdot(\Bvec\Bvec)\right]_i = \partial_j(B_iB_j).
\label{eq:BB-div-components}
\end{equation}
Apply the product rule:
\begin{equation}
\partial_j(B_iB_j)=B_j\partial_j B_i + B_i\partial_j B_j.
\label{eq:BB-div-product}
\end{equation}
Thus,
\begin{equation}
\nabla\cdot(\Bvec\Bvec)=(\Bvec\cdot\nabla)\Bvec+\Bvec(\nabla\cdot\Bvec).
\label{eq:BB-div}
\end{equation}
When $\nabla\cdot\Bvec=0$,
\begin{equation}
\nabla\cdot(\Bvec\Bvec)=(\Bvec\cdot\nabla)\Bvec.
\label{eq:BB-div-simplified}
\end{equation}
Substitute \eqref{eq:isotropic-div} and \eqref{eq:BB-div-simplified} into \eqref{eq:lorentz-pressure-tension}:
\begin{equation}
\Jvec\times\Bvec = -\nabla\cdot\left[\frac12 |\Bvec|^2\Iten - \Bvec\Bvec\right].
\label{eq:lorentz-stress-form}
\end{equation}
The tensor inside the brackets is the magnetic-stress contribution.

\subsection{Final conservative MHD momentum equation}
Insert \eqref{eq:lorentz-stress-form} into \eqref{eq:mhd-mom-with-lorentz}:
\begin{equation}
\pdv{}{t}(\rho\uvec)
+ \nabla\cdot\left[\rho\uvec\uvec + p\Iten + \frac12 |\Bvec|^2\Iten - \Bvec\Bvec\right] = 0.
\label{eq:momentum-expanded}
\end{equation}
Grouping the isotropic terms gives
\begin{equation}
\pdv{}{t}(\rho\uvec)
+ \nabla\cdot\left[\rho\uvec\uvec + \left(p+\frac12 |\Bvec|^2\right)\Iten - \Bvec\Bvec\right] = 0.
\label{eq:momentum-tensor}
\end{equation}
This is the corrected momentum equation used by the solver. The time derivative acts on the \textbf{momentum density} $\rho\uvec$, not on velocity alone.

\subsection{Two-dimensional component form}
The $x$-momentum equation is
\begin{equation}
\pdv{}{t}(\rho u_x)
+ \pdv{}{x}\left(\rho u_x^2 + p + \frac12 |\Bvec|^2 - B_x^2\right)
+ \pdv{}{y}\left(\rho u_xu_y - B_xB_y\right)=0.
\label{eq:xmom2d}
\end{equation}
The diagonal flux term in $x$ contains the isotropic pressure-like part $p+\frac12|\Bvec|^2$ and subtracts $B_x^2$ because magnetic tension opposes compression or stretching along the local field direction.

The $y$-momentum equation is
\begin{equation}
\pdv{}{t}(\rho u_y)
+ \pdv{}{x}\left(\rho u_yu_x - B_yB_x\right)
+ \pdv{}{y}\left(\rho u_y^2 + p + \frac12 |\Bvec|^2 - B_y^2\right)=0.
\label{eq:ymom2d}
\end{equation}
These are exactly the momentum fluxes that later appear in the MATLAB flux functions.

\section{Energy conservation in more detail}
The total energy density is already defined by \eqref{eq:total-energy}. The conservative energy equation used in ideal MHD is
\begin{equation}
\pdv{E}{t} + \nabla\cdot\left[\left(E+\ptot\right)\uvec - (\uvec\cdot\Bvec)\Bvec\right]=0,
\label{eq:energy-cons}
\end{equation}
with
\begin{equation}
\ptot = p + \frac12 |\Bvec|^2.
\label{eq:ptot}
\end{equation}
Here is the physical meaning of the energy flux:
\begin{itemize}[leftmargin=1.8em]
    \item $E\uvec$ is advection of the total stored energy by the bulk flow,
    \item $p\uvec$ is the usual gas-pressure work flux from compressible fluid mechanics,
    \item $\frac12|\Bvec|^2\uvec$ is the advection of magnetic energy,
    \item $-(\uvec\cdot\Bvec)\Bvec$ represents the directional transport associated with the coupling between flow and field.
\end{itemize}
The important practical point for the code is that internal, kinetic, and magnetic energy are not evolved separately. They are packed into a single conserved scalar $E$.

\section{Induction equation in more detail}
The induction equation in ideal MHD is
\begin{equation}
\pdv{\Bvec}{t} - \nabla\times(\uvec\times\Bvec)=0.
\label{eq:induction-curl}
\end{equation}
Its origin is Faraday's law together with the ideal-MHD Ohm's law $\mathbf{E}+\uvec\times\Bvec=0$. Substituting $\mathbf{E}=-\uvec\times\Bvec$ into Faraday's law $\partial \Bvec/\partial t = -\nabla\times \mathbf{E}$ gives \eqref{eq:induction-curl}.

In words, the magnetic field changes because the moving conducting fluid drags and rearranges field lines. In the ideal limit, the field is ``frozen into'' the fluid motion.

\section{Divergence-free constraint}
The magnetic field must also satisfy
\begin{equation}
\nabla\cdot\Bvec = 0.
\label{eq:divB-zero}
\end{equation}
In two dimensions,
\begin{equation}
\pdv{B_x}{x} + \pdv{B_y}{y}=0.
\label{eq:divB-2d}
\end{equation}
This is not a transport equation of the same type as the others; it is a constraint that expresses the absence of magnetic monopoles. Numerically, if $\nabla\cdot\Bvec$ grows too large, the Lorentz-force decomposition above is polluted by unphysical terms. That is why the code monitors $\max |\nabla\cdot\Bvec|$ as a diagnostic.

\section{Flux vectors used by the MATLAB code}
Once the conservative equations are known, the one-dimensional fluxes along each coordinate direction follow directly.

The $x$-direction flux is
\begin{equation}
\Fvec =
\begin{bmatrix}
\rho u_x \\
\rho u_x^2 + \ptot - B_x^2 \\
\rho u_xu_y - B_xB_y \\
\rho u_xu_z - B_xB_z \\
(E+\ptot)u_x - (\uvec\cdot\Bvec)B_x \\
0 \\
u_xB_y-u_yB_x \\
u_xB_z-u_zB_x
\end{bmatrix}.
\label{eq:Fx}
\end{equation}
Each row comes from asking: what amount of the conserved quantity crosses a face whose normal points in $x$?
\begin{itemize}[leftmargin=1.8em]
    \item row 1 is mass flux $\rho u_x$,
    \item rows 2--4 are the three momentum fluxes,
    \item row 5 is the total-energy flux,
    \item rows 6--8 come from the induction equation.
\end{itemize}

The $y$-direction flux is
\begin{equation}
\Gvec =
\begin{bmatrix}
\rho u_y \\
\rho u_yu_x - B_yB_x \\
\rho u_y^2 + \ptot - B_y^2 \\
\rho u_yu_z - B_yB_z \\
(E+\ptot)u_y - (\uvec\cdot\Bvec)B_y \\
u_yB_x-u_xB_y \\
0 \\
u_yB_z-u_zB_y
\end{bmatrix}.
\label{eq:Gy}
\end{equation}
These are the flux vectors passed into the finite-volume and Riemann-solver machinery in the MATLAB implementation.

\section{Summary}
The ideal-MHD equations used by the project come from a clean logical chain:
\begin{enumerate}[leftmargin=1.8em]
    \item derive mass conservation from a control-volume balance,
    \item derive the Euler momentum equation from Newton's second law applied to a fluid element,
    \item convert the Euler equation to conservative form using the continuity equation,
    \item add the Lorentz force density $\Jvec\times\Bvec$,
    \item rewrite the magnetic force as isotropic magnetic pressure plus directional magnetic tension,
    \item rewrite those magnetic terms as the divergence of a magnetic stress tensor,
    \item combine all of the above into conservative fluxes for mass, momentum, energy, and magnetic field.
\end{enumerate}
The point of these notes is that none of the solver equations should look like unexplained magic. Each term has a mechanical origin and a physical interpretation.

\end{document}
